\begin{table}[htbp!]
\centering
\caption{F1-\textit{score} por categoria de consulta e por modelo --- estação de referência (camadas P e N)}
\label{tab:resultados_por_categoria_estacao}
\small
\begin{tabular}{|l|c|c|}
\hline
\textbf{Categoria} & \textbf{Qwen 3 4B} & \textbf{Gemma 4 E2B} \\
\hline
Simples (n=11) & 0,879 & 0,887 \\
Compostas (n=43) & 0,721 & 0,716 \\
Com código MI/INOM (n=15) & 0,886 & 0,867 \\
Com referência temporal relativa (n=15) & 0,760 & 0,655 \\
Com ordenação (n=10) & 0,740 & 0,541 \\
Ambíguas / variações ortográficas (n=35) & 0,694 & 0,693 \\
\hline
\textbf{Média ponderada} & 0,733 & 0,719 \\
\hline
\end{tabular}
\fonte{Elaborado pelo autor. Uma consulta pode pertencer a mais de uma categoria; n = consultas da categoria nas métricas principais. As consultas fora do domínio (categoria F), cujo gabarito é não chamar a ferramenta, não têm campos e aparecem na Figura de acurácia por categoria.}
\end{table}
